\chapter{Thermodynamic Bethe Ansatz}
\label{ch:integrable-thermodynamic-bethe-ansatz}

\ConventionsLorentzian

The thermodynamic Bethe ansatz turns the exact two-body scattering data of a
massive integrable theory into finite-temperature or finite-size
thermodynamics.  The construction in this chapter is stated first for a
diagonal elastic \(S\)-matrix.  Matrix-valued scattering, bound-state strings,
and nested Bethe ansatz data require additional species and auxiliary-root
bookkeeping but obey the same variational logic once the allowed densities
are specified.

\section{Bethe--Yang Quantization}

Place the theory on a spatial circle of circumference \(L\).  For \(n\)
particles of masses \(m_a\), rapidities \(\theta_j\), and diagonal scattering
phases
\[
  S_{ab}(\theta)=\exp\!\bigl(\ii\delta_{ab}(\theta)\bigr),
\]
periodicity of the wavefunction of particle \(j\) gives the Bethe--Yang
equation
\[
  m_{a_j}L\sinh\theta_j
  +\sum_{k\neq j}\delta_{a_ja_k}(\theta_j-\theta_k)
  =
  2\pi I_j,
  \qquad I_j\in\mathbb Z+\nu_j .
\]
The shift \(\nu_j\) depends on statistics and boundary convention.  The
one-particle energy and momentum are
\[
  E_a(\theta)=m_a\cosh\theta,\qquad
  p_a(\theta)=m_a\sinh\theta .
\]

\section{Density Equations}

In the thermodynamic limit \(L\to\infty\), introduce particle and hole
densities \(\rho_a(\theta)\) and \(\rho_a^h(\theta)\), normalized so that
\[
  L\rho_a(\theta)d\theta
\]
is the number of occupied rapidities of species \(a\) in \(d\theta\).  Define
the scattering kernel
\[
  \varphi_{ab}(\theta)
  =
  \frac{1}{2\pi}\frac{d}{d\theta}\delta_{ab}(\theta)
  =
  \frac{1}{2\pi\ii}\frac{d}{d\theta}\log S_{ab}(\theta).
\]
Differentiating the Bethe--Yang equations gives the density constraint
\[
  \rho_a(\theta)+\rho_a^h(\theta)
  =
  \frac{m_a}{2\pi}\cosh\theta
  +\sum_b\int_{\mathbb R}d\theta'\,
  \varphi_{ab}(\theta-\theta')\rho_b(\theta').
\]
This equation is purely kinematic: it counts available rapidity levels after
the scattering phase shifts have shifted the quantization conditions.

\section{Entropy and Free Energy}

For fixed total density \(\rho_a^t=\rho_a+\rho_a^h\), the number of
microscopic occupancy choices in an interval is
\[
  \binom{L\rho_a^t\,d\theta}{L\rho_a\,d\theta}.
\]
Stirling's formula gives the entropy density functional
\[
  s[\rho,\rho^h]
  =
  \sum_a\int d\theta\,
  \left[
  \rho_a^t\log\rho_a^t
  -\rho_a\log\rho_a
  -\rho_a^h\log\rho_a^h
  \right].
\]
At inverse temperature \(R\), the free energy per length is
\[
  f[\rho,\rho^h]
  =
  \sum_a\int d\theta\,m_a\cosh\theta\,\rho_a(\theta)
  -\frac1R s[\rho,\rho^h],
\]
subject to the density constraint.

\section{Pseudoenergy Equation}

Define the pseudoenergy by
\[
  \epsilon_a(\theta)=\log\frac{\rho_a^h(\theta)}{\rho_a(\theta)} .
\]
Varying the constrained free energy with respect to \(\rho_a\) gives
\[
  \epsilon_a(\theta)
  =
  Rm_a\cosh\theta
  -\sum_b\int_{\mathbb R}d\theta'\,
  \varphi_{ab}(\theta-\theta')
  \log\!\left(1+e^{-\epsilon_b(\theta')}\right).
\]
This is the thermodynamic Bethe ansatz equation.  The equilibrium free energy
is
\[
  f(R)
  =
  -\frac{1}{2\pi R}
  \sum_a m_a\int_{\mathbb R}d\theta\,
  \cosh\theta\,
  \log\!\left(1+e^{-\epsilon_a(\theta)}\right).
\]

\begin{proposition}[Variational derivation of TBA]
\label{prop:tba-variational-derivation}
For diagonal elastic scattering with smooth phase kernels and a thermodynamic
limit in which the Bethe--Yang root densities exist, the stationary point of
the constrained free-energy functional is exactly the pseudoenergy equation
above.
\end{proposition}

\begin{proof}
Let \(\rho_a^t=\rho_a+\rho_a^h\).  A variation of \(\rho_a\) changes
\(\rho_c^t\) through the density constraint by
\[
  \delta\rho_c^t(\theta)
  =
  \sum_a\int d\theta'\,
  \varphi_{ca}(\theta-\theta')\delta\rho_a(\theta').
\]
The entropy variation is
\[
  \delta s
  =
  \sum_a\int d\theta\,
  \left[
  \log\frac{\rho_a^h}{\rho_a}\,\delta\rho_a
  +
  \log\frac{\rho_a^t}{\rho_a^h}\,\delta\rho_a^t
  \right].
\]
Using \(\log(\rho_a^t/\rho_a^h)=\log(1+e^{-\epsilon_a})\), substituting the
constraint variation, and setting \(\delta f=0\) for arbitrary
\(\delta\rho_a\) gives the displayed integral equation.  Inserting the
stationary relation back into \(f\) gives the stated free energy.
\end{proof}

\section{Finite-Size Ground-State Energy}

Interchanging space and Euclidean time maps the thermal free energy on an
infinite line at inverse temperature \(R\) to the ground-state energy on a
circle of circumference \(R\), after the mirror-channel scattering data have
been identified.  In a relativistic theory the mirror dispersion is obtained
by a rapidity shift by \(\ii\pi/2\), so the same TBA equation gives
\[
  E_0(R)
  =
  -\frac{1}{2\pi}
  \sum_a m_a\int_{\mathbb R}d\theta\,
  \cosh\theta\,
  \log\!\left(1+e^{-\epsilon_a(\theta)}\right).
\]
The effective central charge extracted from the ultraviolet limit is
\[
  c_{\rm eff}(R)
  =
  -\frac{6R}{\pi}E_0(R),
  \qquad R\to0,
\]
provided the limit exists and the infrared bulk energy has been subtracted in
the declared normalization.

\section{Free Majorana Example and the Ultraviolet Central Charge}

The massive Ising field theory in its fermionic phase gives the minimal
diagonal example.  It has one neutral particle of mass \(m\).  The two-particle
scattering phase is the constant
\[
  S(\theta)=-1 .
\]
The derivative kernel is therefore zero:
\[
  \varphi(\theta)
  =
  \frac{1}{2\pi\ii}\frac{\dd}{\dd\theta}\log S(\theta)=0,
\]
after the constant branch of \(\log(-1)=\ii\pi\) has been fixed.  The TBA
equation reduces to
\[
  \epsilon(\theta)=r\cosh\theta,
  \qquad r=mR .
\]
The finite-size ground-state energy in the normalization of the preceding
section is
\begin{equation}
  E_0(R)
  =
  -\frac{m}{2\pi}
  \int_{-\infty}^{\infty}\dd\theta\,
  \cosh\theta\,
  \log\!\left(1+\ee^{-r\cosh\theta}\right).
  \label{eq:tba-free-majorana-ground-energy}
\end{equation}
The ultraviolet central charge follows from a direct scaling limit.  Put
\[
  x=\sinh\theta,\qquad \dd x=\cosh\theta\,\dd\theta .
\]
Then
\[
  I(r)
  :=
  \int_{-\infty}^{\infty}\dd\theta\,
  \cosh\theta\,
  \log(1+\ee^{-r\cosh\theta})
  =
  \int_{-\infty}^{\infty}\dd x\,
  \log(1+\ee^{-r\sqrt{1+x^2}}).
\]
Rescale \(y=rx\).  This gives
\[
  rI(r)
  =
  \int_{-\infty}^{\infty}\dd y\,
  \log(1+\ee^{-\sqrt{r^2+y^2}}).
\]
For \(0<r\leq1\) the integrand is bounded by
\(\log(1+\ee^{-|y|})\), which is integrable on \(\mathbb R\), and it
converges pointwise to \(\log(1+\ee^{-|y|})\).  Dominated convergence gives
\[
  \lim_{r\downarrow0}rI(r)
  =
  2\int_0^\infty\dd y\,\log(1+\ee^{-y}).
\]
Expanding \(\log(1+\ee^{-y})\) for \(y>0\) and integrating termwise gives
\[
  \sum_{k=1}^{\infty}\frac{(-1)^{k+1}}{k^2}
  =
  \int_0^\infty\dd y\,\log(1+\ee^{-y})
  =
  \frac{\pi^2}{12},
\]
Equation~\eqref{eq:tba-free-majorana-ground-energy} gives
\[
  E_0(R)
  =
  -\frac{\pi}{12R}+o(R^{-1}).
\]
Thus
\[
  c_{\rm eff}
  =
  \lim_{R\downarrow0}\left(-\frac{6R}{\pi}E_0(R)\right)
  =
  \frac12 .
\]
The result matches the central charge of the ultraviolet Majorana conformal
field theory, but the computation above uses only the declared finite-size
TBA datum and the normalization of \(E_0(R)\).

\section{Theorem Boundary}

The derivation above assumes existence of thermodynamic root densities,
validity of Bethe--Yang quantization up to errors negligible at large \(L\),
and uniqueness of the free-energy minimizer.  In concrete models these are
mathematical hypotheses to be proved or checked.  Bound states, non-diagonal
scattering, zero modes, and massless flows change the density system and must
be included explicitly.

\begin{openproblem}[TBA from local QFT]
\label{op:tba-from-local-qft}
Derive the TBA equations directly from a constructed local QFT with
factorized scattering: prove Bethe--Yang quantization with controlled
finite-volume errors, establish thermodynamic root-density convergence, and
identify the finite-size ground-state energy with the mirror free energy.
\end{openproblem}
